\section{Open questions raised by this replication}

Five concrete open questions grounded in the paper's declared scope and left un-tested by this replication. Each has a stated basis and a concrete next-step probe.

\begin{enumerate}
  \item \textbf{Relativistic / electromagnetic extension of the energy-conserving split.}
  Does the Eq.~5 field-parallel velocity correction remain exactly conservative in a relativistic/EM PIC context where the velocity update is Lorentz-factor coupled and $q\Phi$ is replaced by the full four-potential contribution?
  \emph{Basis:} The paper derives the correction in an electrostatic, non-relativistic setting inside a single element with constant $\mathbf{E}$. Aleph as described is electrostatic; the relativistic/EM extension is not addressed.
  \emph{Next step:} Re-derive the split-child velocity equation replacing $\tfrac{1}{2}mv^2$ with $(\gamma-1)mc^2$ and $q\Phi$ with $q(\Phi - \mathbf{v}\cdot\mathbf{A}/c)$; implement in a stand-alone Boris-pusher test with a prescribed EM field gradient; measure relative energy error across the same 20{,}000-split sweep at $\gamma\in\{1.0, 1.1, 2.0, 10.0\}$ to quantify any leak.

  \item \textbf{Interaction with adaptive mesh refinement (AMR).}
  Does the element-local field assumption (constant $\mathbf{E}$ within an element in the Eq.~5 correction) still hold when children land in a newly-refined sub-element with a different local $\mathbf{E}$? Does inter-level migration preserve the invariants?
  \emph{Basis:} Sec.~3.1 uses a nearest-face bound $\mathrm{d}x_{\max}$ and assumes constant $\mathbf{E}$ within the parent's element. Aleph as tested is uniform-mesh.
  \emph{Next step:} Build a two-level nested mesh test where a parent macroparticle sits near a refinement boundary; sweep child-displacement magnitudes so a fraction of children cross into the refined level; measure per-event energy error and count boundary-straddling events.

  \item \textbf{GPU / exascale scaling of the operators.}
  Does the per-cell serialization implied by the pairwise merge (Eqs.~6--12 need a nearest-speed pairing inside each cell) become the throughput bottleneck at the Aleph scales targeted by Sandia/NVIDIA?
  \emph{Basis:} The author list includes NVIDIA and the paper claims performance, but the paper focuses on physics correctness rather than a GPU strong/weak-scaling study. Aleph internals are not public.
  \emph{Next step:} Prototype the merge kernel in CUDA or SYCL with three pairing strategies (per-cell nearest-speed serial, atomic-buffered pairwise, and approximate bucketed pairing); benchmark on a single H100 at particle counts $10^6$--$10^8$ per rank; compare against the split kernel throughput to isolate the true bottleneck.

  \item \textbf{Sensitivity to cutoffs and weight window.}
  How sensitive are the physical outputs (VDF shape, ionization rate, wall heat flux) to the merge cutoffs (Eq.~9: $\Delta v < v_{\mathrm{th}}$, $\alpha < 30^\circ$) and to the target-weight window $(W_{\min}, W_{\max})$?
  \emph{Basis:} The paper fixes cutoffs at $v_{\mathrm{th}}(T{=}5\,\mathrm{eV})$ and $30^\circ$ without a sensitivity sweep. Our replication recovered KE-loss statistics only at these fixed values.
  \emph{Next step:} In the 0D swarm test, sweep the angular cutoff over $\{10, 20, 30, 45, 60\}^\circ$ and the speed cutoff over $\{0.5, 1.0, 2.0\}\times v_{\mathrm{th}}$; measure fitted $\beta$, EEDF second moment, and cumulative merge-KE loss; identify the largest cutoff window that keeps thermal-energy drift under 1\% per $e$-fold of growth.

  \item \textbf{Head-to-head against alternative coalescence schemes.}
  How does this split/merge scheme compare quantitatively against alternatives (Lapenta--Brackbill $v$-space Voronoi merging, Assous $k$-means merging, Welch--Rose weighted-density resampling) on the same 0D H$_2$ growth test and Test-4.1 sheath problem?
  \emph{Basis:} The paper motivates the method as an improvement over prior schemes but does not provide a head-to-head numerical benchmark against a named alternative. Our replication only exercised the paper's own scheme.
  \emph{Next step:} Implement one alternative (Lapenta--Brackbill Voronoi merge is the smallest lift) in the same reimplementation harness; run both schemes on the identical 0D growth test at $N_c\in\{100, 1000, 10000\}$; report side-by-side fitted-$\beta$ bias, computational-count bounds, cumulative KE drift, and per-event cost.
\end{enumerate}
